% Capítulo 3: Necesidades de monitoreo.
\chapter{Necesidades relacionadas con monitoreo de temperatura y humedad relativa
en la detección de \textit{Tetranychus urticae}}
\label{cap:necesidades-monitoreo}

El contexto de estudio corresponde a plantas madre de cítricos cultivadas bajo
cubierta en viveros de la zona rural de Popayán, Cauca. Su importancia para la
propagación del material vegetal motiva la caracterización del entorno y de las
condiciones de vigilancia fitosanitaria \cite{ica2019resolucion}. Esta
caracterización constituye el primer paquete de trabajo previsto en el anteproyecto.

La temperatura y la humedad relativa son variables de interés para interpretar
las condiciones asociadas al desarrollo de \textit{Tetranychus urticae}
\cite{praslicka2004}. Sus registros aportarán contexto al sistema cognitivo, pero
no se considerarán por sí solos una confirmación de presencia o ausencia de la
plaga. La detección directa deberá contrastarse con observaciones de la planta y
una referencia de identificación experta.

Se prevé realizar observaciones de campo y entrevistas con expertos para
caracterizar el hábitat de la plaga, sus efectos visibles y las prácticas de
inspección. A partir de esta información se definirán los requisitos de
adquisición ambiental y visual, la ubicación de los puntos de monitoreo y la
representación del conocimiento agrícola. Las restricciones de conectividad,
energía y acceso a las plantas deberán documentarse en el sitio de estudio.

\textit{[Pendiente: incorporar los registros de caracterización, los requisitos
validados y la justificación de rangos, precisión y frecuencia de muestreo.
Las actividades anteriores describen el plan, no visitas o mediciones realizadas.]}
